 \documentclass[11pt]{amsart}
\usepackage[utf8]{inputenc}
\usepackage{bm}
\usepackage{fontenc}
\usepackage{amsfonts}
\usepackage{amssymb}
\usepackage{amsmath}
\usepackage{amsthm}\usepackage{mathtools}
\usepackage{enumerate}
\usepackage{enumitem} 
\usepackage[pagebackref,colorlinks,linkcolor=blue,citecolor=blue,urlcolor=blue,hypertexnames=false]{hyperref}
\usepackage{mathrsfs}
\usepackage{tikz}
\usetikzlibrary{calc}
\usepackage{marginnote}
\usepackage{xcolor,enumitem}
\usepackage{soul}
\usepackage[all,cmtip]{xy} \usepackage{caption} 
	
 
 

 
 
\newcommand{\N}{\mathbb{N}}
 
 

 
\newcommand{\Lip}{\mathrm{Lip}}
  
  

 

%\numberwithin{equation}{section}

\newtheorem{theorem}{Theorem} 
\newtheorem*{theorem*}{Theorem}
\newtheorem{proposition}[theorem]{Proposition}
\newtheorem{problem}[theorem]{Problem}
\newtheorem*{proposition*}{Proposition}
\newtheorem{lemma}[theorem]{Lemma}
\newtheorem*{lemma*}{Lemma}
\newtheorem{corollary}[theorem]{Corollary}
\newtheorem*{corollary*}{Corollary}
\newtheorem{fact}[theorem]{Fact}
\newtheorem*{fact*}{Fact}
\theoremstyle{definition}
\newtheorem{definition}[theorem]{Definition}
\newtheorem*{definition*}{Definition}
\newtheorem*{Case1}{\textbf{Case 1}}
\newtheorem*{Case2}{\textbf{Case 2}} 

\newtheorem*{acknowledgements*}{Acknowledgements}
\newtheorem*{disclaimer*}{Disclaimer on originality of the work}
\newtheorem*{leanformalization*}{Lean formalization}
\newtheorem*{AIusage}{AI usage statement}
\newtheorem{claim}[theorem]{Claim}
\newtheorem*{claim*}{Claim}
\newtheorem*{proofsketch}{{\normalfont\textit{Sketch of the proof}}}
\newtheorem*{subclaim}{Subclaim}
\newtheorem{conjecture}[theorem]{Conjecture}
\newtheorem*{conjecture*}{Conjecture}
\newtheorem{notation}{Notation}
\newtheorem{theoremi}{Theorem}
\newtheorem{assumption}[theorem]{Assumption}
\newtheorem{convention}[theorem]{Convention}
\newtheorem{question}[theorem]{Question}

\renewcommand\thetheoremi{\Alph{theoremi}}

\theoremstyle{remark}
\newtheorem{example}[theorem]{Example}
\newtheorem*{example*}{Example}
\newtheorem{remark}[theorem]{Remark}
\newtheorem*{remark*}{Remark}
\newtheorem{note}{Note}
\newtheorem*{note*}{Note}
%\newtheorem{question}{Question}
\newtheorem*{question*}{Question}
\newtheorem{exmpl}[theorem]{Example}

 
\newcommand{\norm}[1]{\left\lVert #1 \right\rVert}
 
 
\DeclareMathOperator{\supp}{supp}
 
\DeclareMathOperator{\diam}{diam}
 
 

 \newcommand{\cg}[1]{{\color{blue}#1}}

 
  
 \usepackage{anysize}
\usepackage{float}

\usepackage{anysize}
\marginsize{2.9cm}{2.9cm}{2.9cm}{2.9cm}
 


\begin{document}

\date{\today} % Activate to display a given date or no date

\title[Expanders prevent Property (H)]{Expanders prevent Property (H)}
 
  
\author[Braga]{Bruno M. Braga}
\address[B. M. Braga]{IMPA, Estrada Dona Castorina 110, 22460-320, Rio de Janeiro, Brazil}
\email{demendoncabraga@gmail.com}
\urladdr{\url{https://sites.google.com/site/demendoncabraga}}
\thanks {B. M. Braga  was partially supported by FAPERJ, grant E-26/200.167/2023,  by CNPq, grant 303571/2022-5, and by Serrapilheira, grant R-2501-51476.}

\author[Gartland]{C. Gartland}
\address[C. Gartland]{Department of Mathematics and Statistics, University of
North Carolina at Charlotte, Charlotte, North Carolina, USA}
\email{cgartla1@charlotte.edu}
\urladdr{\url{https://chrisgartland.wordpress.com/}}
\thanks{C. Gartland was partially supported by NSF Award DMS-2555144.}

\author[Lancien]{G. Lancien}
\address[G. Lancien]{Universit\'e Marie et Louis Pasteur, CNRS, LmB (UMR 6623), F-25000 Besan\c con, France.}
\email{gilles.lancien@univ-fcomte.fr}
%\urladdr{\url{ }}
\thanks{G. Lancien was partially supported by the French ANR project No. ANR-24-CE40-0892-01.}

\author[Motakis]{P. Motakis}
\address[P. Motakis]{Department of Mathematics and Statistics, York University, 4700 Keele Street, Toronto, Ontario, M3J 1P3, Canada}
\email{pmotakis@yorku.ca}
\urladdr{\url{https://pmotakis.mathstats.yorku.ca/}}
\thanks{P. Motakis was partially supported by NSERC Grant RGPIN-2021-03639.}

\author[Perneck\'a]{E. Perneck\'a}
\address[E. Perneck\'a]{Faculty of Information Technology, Czech Technical University in Prague, Th\'akurova 9, 160 00, Prague 6, Czech Republic}
\email{eva.pernecka@fit.cvut.cz}
%\urladdr{\url{ }}
\thanks{}

\author[Schlumprecht]{Th. Schlumprecht}
\address[Th. Schlumprecht]{Texas A\&M University, College Station, TX 77843,
USA and Faculty of Electrical Engineering, Czech Technical University in Prague, Zikova 4, 166 27, Prague}
\email{t-schlumprecht@tamu.edu}
\urladdr{\url{https://people.tamu.edu/~t-schlumprecht/}}
\thanks{Th. Schlumprecht was partially supported by NSF Award DMS-2349322.}
\begin{abstract}
    We show that if a sequence of expander graphs equi-coarsely embeds into a Banach space, then this Banach space  fails Kasparov and Yu's Property (H).  We provide a Lean verification of this statements.
\end{abstract}

\maketitle


Kasparov and Yu have introduced in \cite{KasparovYu2012GeoTop} a property for Banach spaces called \emph{Property (H)} and showed that the Novikov conjecture holds for any countable group which coarsely embeds into a Banach space satisfying it (\cite[Theorem 1.2]{KasparovYu2012GeoTop}). While many examples of Banach spaces with Property (H) can be found by using Mazur's modified maps constructed by Odell and Schlumprecht in \cite{OdellSchlumprecht1994Acta} (e.g., this is the case for all separable Banach lattices with nontrivial cotype, see \cite[Theorem 1.1]{ChengWang2018JMAA}), there were no known examples of separable Banach spaces failing Property (H). In particular, it was open whether $c_0$ has Property (H); this was seen  as a special test case since $c_0$ is Lipschitz universal for all separable metric spaces (see \cite[Theorem on Page 288]{Aharoni1974Israel}). This short note shows that not only does $c_0$ not have Property (H), but in fact no Banach space containing expander graphs coarsely has Property (H). This is an immediate consequence of our main result, which reads  as follows (throughout, $S_X$ denotes the unit sphere of the Banach space $X$).




\begin{theorem}\label{Thm.main}
    Let $(X_n)_n$ and $(Y_n)_n$ be sequences of Banach spaces such that $Y_n\subseteq L_1$ for all $n\in\N$. Let $(G_n=(V_n,E_n))_n$ be a sequence of expander graphs and assume that there is a sequence of equi-coarse embeddings $(V_n\to X_n)_n$. Let $(F_n\colon S_{X_n}\to S_{Y_n})_n$ be a sequence of equi-uniformly continuous maps. Then, for $n\in\N$ large enough, the maps $F_n$ are all null-homotopic.
\end{theorem}


We briefly remind the reader of the terminologies in Theorem \ref{Thm.main}. Firstly, given $k\in\N$ and $h>0$, a finite   graph $G=(V,E)$, where $V$ is the set of \emph{vertices} and $E$ the set of \emph{edges} of $G$, is called a \emph{$(k,h)$-expander} if each vertex has at most $k$ neighbors and 
\[|\{v\in V\setminus A\mid \exists u\in A, \ (v,u)\in E\}|\geq h|A|\]
for all $A\subseteq V$ with $|A|\leq |V|/2$. A \emph{sequence of expander graphs} is understood to be a sequence of $(k,h)$-expanders, for some $k\in\N$ and $h>0$, such that $\lim_n|V_n|=\infty$. As expanders must be connected, we view the vertex set of an expander as a metric space by endowing it with the shortest path distance, which we denote by $d_n$.  For  more on expander graphs, see the monograph \cite{LubotzkyBook2010}. We then say that   $(\varphi_n\colon V_n\to X_n)_n$ is a sequence of  \emph{equi-coarse embeddings} if there are $L>0$ and a nondecreasing $\rho\colon [0,\infty)\to[0,\infty)$ with $\lim_{t\to \infty}\rho(t)=\infty$  such that 
\[\rho(d_n(v,u))\leq \|\varphi_n(v)-\varphi_n(u)\|\leq Ld_n(v,u)\ \text{ for all }n\in\N\ \text{ and } v,u\in V_n.\]



In light of Theorem \ref{Thm.main}, in order to rule out Property (H) for Banach spaces containing expanders, it is enough to remind the reader of the definition of this property. Indeed, a Banach space $X$ has \emph{Property (H)} if there are a uniformly continuous function $F\colon S_X\to S_{\ell_2}$ and increasing sequences of finite dimensional Banach spaces $(X_n)_n$ and $(H_n)_n$ of $X$ and of the Hilbert space $\ell_2$, respectively, such that $X=\overline{\bigcup_nX_n}$,  each $F(S_{X_n})\subseteq S_{H_n}$, and each $F_{\restriction S_{X_n}}\colon  S_{X_n}\to S_{H_n}$ is a degree one map.  As a degree one map between finite dimensional spheres is never null-homotopic, the following is an immediate consequence of Theorem \ref{Thm.main}.

\begin{corollary}\label{Cor.Prop.H}
If a sequence of expander graphs equi-coarsely embeds into a Banach space, then the Banach space  does not have Property (H). In particular, if a Banach space contains the finite dimensional $(\ell_\infty^n)_n$ with uniform distortion (such as $c_0$), then it cannot have Property (H).\qed
\end{corollary}

 We remark that it is currently unknown if a Banach space which does not  contain the finite dimensional $(\ell_\infty^n)_n$ with uniform distortion  can equi-biLipschitzly contain a sequence of expander graphs.

As there are finitely generated groups whose Cayley graphs contain expanders isometrically (see \cite[Theorem 4]{Osajda2020Acta}), the following corollary is also immediate.

\begin{corollary}\label{Cor.no.PropH.cL.universal}
    There is no Banach space with Property (H) which contains coarse  copies of all countable groups.\qed
\end{corollary} 


Another consequence of Theorem \ref{Thm.main} is a negative answer to a question of W. B. Johnson posted in \cite{JohnsonOverFlow}. Precisely, since any metric space with $n$ elements isometrically embeds into $\ell_\infty^n$,  Theorem \ref{Thm.main} also has the following corollary. 

\begin{corollary}\label{Corollary.Johnson} 
    There is no sequence of homeomorphisms $(F_n\colon S_{\ell_\infty^n}\to S_{\ell_2^n})_n$ which is equi-uniformly continuous.\qed
\end{corollary}

We now start presenting the proof of Theorem \ref{Thm.main}. The arguments are  surprisingly simple and  only make use of a   standard Poincaré inequality satisfied by expanders: If $G=(V,E)$ is a $(k,h)$-expander, then  for any $f\colon V\to L_1$ we have
     \begin{equation}\label{PIII}
\frac{1}{|V|}\sum_{v\in V}\left\|f(v)-\frac{1}{|V|}\sum_{u\in V}f(u)\right\|\leq \frac{2k}{h}\Lip(f) \tag{\textbf{P.I.}}
  \end{equation}
     (see \cite[Theorem 4.7]{Ostrovskii2013Book}).
     
 Given a Banach space $X$, $B_X$ denotes its closed unit ball. We start with a lemma which modifies the equi-coarse embeddings $(V_n\to X_n)_n$ in Theorem \ref{Thm.main} to maps which better suit our purposes. 

\begin{lemma}\label{Lemma}
    Let $(X_n)_n$ be Banach spaces and $(G_n=(V_n,E_n))_n$ be a sequence of expander graphs. If there is a sequence of  equi-coarse embeddings $(V_n\to X_n)_n$, then there is a sequence of functions $(\psi_n\colon V_n\to B_{X_n})_n$ such that 
    \begin{equation}\label{Eq.prop.psi}
    \sum_{v\in V_n}\psi_n(v)=0, \ \liminf_{n\to \infty}\frac{1}{|V_n|}\sum_{v\in V_n}\|\psi_n(v)\|\geq 1\ \text{ and } \lim_{n\to\infty}\Lip(\psi_n)=0.
    \end{equation}
\end{lemma}

\begin{proof} 
Fix $k\in\N$ and $h>0$ such that each $G_n$ is a $(k,h)$-expander. For each $n\in\N$, let $d_n$ be the shortest path metric on $V_n$ and let $(\varphi_n\colon V_n\to {X_n})_n$ be a sequence of equi-coarse embeddings; so there is $L>0$ and $\rho\colon [0,\infty)\to [0,\infty)$ with $\lim_{t\to \infty}\rho(t)=\infty$ such that 
\[\rho(d_n(v,u))\leq \|\varphi_n(v)-\varphi_n(u)\|\leq Ld_n(v,u) \ \text{ for all }\ n\in\N\text{ and }\ v,u\in V_n.\]
By translating the $\varphi_n$'s  if needed,  assume that $\sum_{v\in V_n}\varphi_n(v)=0$.
We start noticing that 
\begin{equation}\label{Eq.Antoinfty}A_n=\frac{1}{|V_n|}\sum_{v\in V_n}\|\varphi_n(v)\|\underset{n\to \infty}{\longrightarrow} \infty.\end{equation}
Indeed, as each $G_n$ is a $(k,h)$-expander, the condition on $k$ implies that for each $r>0$, the balls of radius $r$ of each $V_n$ contain at most $k^{r+1}$ elements. Therefore, 
\[2A_n\geq \frac{1}{|V_n|^2} \sum_{v,u\in V_n}\|\varphi_n(v)-\varphi_n(u)\|\geq \left(1-\frac{k^{r+1}}{|V_n|}\right)\rho(r).\]
Letting $n\to \infty$ and then $r\to \infty$, we obtain that  $\lim_{n\to\infty}A_n=\infty$.


The proof now  consists of truncating the functions $\varphi_n$ so that they lie in the balls $A_nB_{X_n}$, translating them so that their mean is zero, and then rescaling them appropriately. For each $n\in\N$, let $R_n\colon X_n\to A_nB_{X_n}$ denote the radial truncation
\begin{equation*}
    R_n(x)=\left\{\begin{array}{cc}
      x   ,&  \text{ if }\ \|x\|\leq A_n,\\
       A_n\frac{x}{\|x\|},  & \text{ if }\ \|x\|>A_n. 
    \end{array}\right.
\end{equation*}
It is straightforward to check that $\Lip(R_n)\leq 2$. For $n\in\N$ and $v\in V_n$, define 
\[\theta_n(v)=R_n(\varphi_n(v))\ \text{ and }\ m_n=\frac{1}{|V_n|}\sum_{u\in V_n}\theta_n(u).\]
Let   $C=2kL/h$ and define each $\psi_n\colon V_n\to B_{X_n}$ by letting \[\psi_n(v)=\frac{\theta_n(v)-m_n}{A_n+C}\   \text{ for all } \ v\in V_n.\]
Let us now show that these maps have the required properties. The fact that $\sum_{v\in V_n}\psi_n(v)=0$ is immediate and, as $\lim_{n\to \infty}A_n=\infty$, we also have that 
\[\Lip(\psi_n)\leq \frac{2L}{A_n+C}\underset{n\to \infty}{\longrightarrow} 0.\]


We are left to show that each $\psi_n$ takes values in $B_{X_n}$ and that the limit inferior of  $\tfrac{1}{|V_n|}\sum_{v\in V_n}\|\psi_n(v)\|$  is at least $1$.  As each $v\in V_n\mapsto \|\varphi_n(v)\|\in \mathbb{R}$ has Lipschitz constant at most $L$, \eqref{PIII} and the choice of $C$ give that 
\begin{equation}\label{Eq.piii}
    \frac{1}{|V_n|}\sum_{v\in V_n}\left|\|\varphi_n(v)\|-A_n\right|\leq C.
\end{equation}
Since $\sum_{v\in V_n}\varphi_n(v)=0$, \eqref{Eq.piii} implies that
\begin{equation*}
    \|m_n\|\leq\frac{1}{|V_n|}\sum_{v\in V_n}\|\varphi_n(v)-\theta_n(v)\|\leq \frac{1}{|V_n|}\sum_{v\in V_n}|\|\varphi_n(v)\|-A_n|\leq C.
\end{equation*}
By the formula of $\psi_n$, this  implies that these maps take values in $B_{X_n}$ as needed.
Also, since 
\[\|\theta_n(v)\|=\min\{\|\varphi_n(v)\|,A_n\}\geq \|\varphi_n(v)\|-|\|\varphi_n(v)\|-A_n|,\]
 \eqref{Eq.piii}  implies that
\begin{equation*}
    \frac{1}{|V_n|}\sum_{v\in V_n}\|\theta_n(v)\|\geq     \frac{1}{|V_n|}\sum_{v\in V_n}\|\varphi_n(v)\|-    \frac{1}{|V_n|}\sum_{v\in V_n}|\|\varphi_n(v)\|-A_n|\geq A_n-C.
\end{equation*}
We then conclude that 
\begin{align*}
     \frac{1}{|V_n|}\sum_{v\in V_n}\|\psi_n(v)\|\geq \frac{\frac{1}{|V_n|}\sum_{v\in V_n}\|\theta_n(v)\|-C}{A_n+C}\geq \frac{A_n-2C}{A_n+C}\underset{n\to \infty}{\longrightarrow}1,
\end{align*}
as desired.
\end{proof}

 
 
 
 Another trivial observation needed for the proof of Theorem \ref{Thm.main} is the following sufficient condition for null-homotopy of our maps:

 \begin{proposition}\label{Prop.null.homotopy}
     Let $X$ and $Y$ be Banach spaces and $f\colon B_X\to B_Y$ be a continuous map such that $\inf_{x\in B_X}\|f(x)\|>0$. Then the map $F\colon S_X\to S_Y$ given by $F(x)=f(x)/\|f(x)\|$ is null-homotopic.
 \end{proposition}
 
 \begin{proof}
Let $g_s\colon B_X\to B_X$, $s\in [0,1]$, be a homotopy from $g_0=\mathrm{Id}_{B_X}$ to the constant map $g_1\equiv 0$. Then the map $G_s(x)=f(g_s(x))/\|f(g_s(x))\|$ defines a homotopy \cg{in $S_Y$} from $G_0=F$ to the constant map $G_1\equiv f(0)/\|f(0)\|$.    
\end{proof}
 

 

\begin{proof}[Proof of Theorem \ref{Thm.main}]
We start fixing the objects given in the statement of the theorem together with Lemma \ref{Lemma}. Discarding finitely many terms, we may assume that $|V_n|\geq5$ for all $n\in\N$. Let  $k\in\N$ and $h>0$ be such that each $G_n=(V_n,E_n)$ is a $(k,h)$-expander. For each $n\in\N$ let $d_n$ be the shortest path metric on $V_n$. Let $(\psi_n\colon V_n\to B_{X_n})_n$ be the maps given by Lemma \ref{Lemma} due to the existence of  equi-coarse embeddings $(V_n\to X_n)_n$ and recall that  $(F_n\colon S_{X_n}\to S_{Y_n})_n$ denotes a  sequence of equi-uniformly continuous maps.
 


We  now  show how to use Proposition \ref{Prop.null.homotopy} in   order to obtain the null-homotopy of $F_n$ for large $n$. The proof has two parts: (1) we first show each $F_n$ to be homotopic to some other map, and then (2) we show that this other map is the normalization of a function  on $B_{X_n}$ which avoids zero. We start extending each $F_n$ radially to $G_n:2B_{X_n}\to2B_{Y_n}$ by letting \[
  G_n(x)=
  \begin{cases}
    \norm{x}F_n\left(\frac{x}{\norm{x}}\right),&x\neq0,\\
    0,&x=0.
  \end{cases}
\]
Clearly,   each $G_n$ is  norm preserving and its restriction to $S_{X_n}$ is $F_n$. It is also straightforward to check that, as $(F_n)_n$ are equi-uniformly continuous,  so are $(G_n)_n$ (e.g., see \cite[Proposition 2.9]{OdellSchlumprecht1994Acta}). Fix then an increasing $\omega\colon [0,\infty)\to [0,\infty)$  with $\lim_{t\to 0}\omega(t)=0$ such that
\[\|G_n(x)-G_n(y)\|\leq \omega(\|x-y\|)\ \text{ for all }\ n\in\N\ \text{ and  }x,y\in 2B_{X_n}.\] 
 
 

 
  For each  $n\in\N$,   $ 
  x\in B_{X_n}$, and  $s\in[0,1]$,  let 
\begin{equation*}
 H^n_s(x)=\frac{1}{|V_n|}\sum_{v\in V_n} G_n(x+s\psi_n(v)).
  \label{eq:average}
\end{equation*}
Then, for each  $x\in B_{X_n}$ and $s\in[0,1]$, applying \eqref{PIII} to the map   
\[
v\in V_n\mapsto G_n(x+s\psi_n(v))\in Y_n\subseteq  L_1 ,
\]
we obtain that, for all $n\in\N$,
\begin{equation*}
  \begin{aligned}
   \frac{1}{|V_n|} \sum_{v\in V_n} \norm{G_n(x+s\psi_n(v))-H^n_s(x)}
    &\leq \frac{2k}{h}\omega\left(\Lip(\psi_n)\right).
  \end{aligned}
  \label{eq:deviation}
\end{equation*}
As each $G_n$ is norm preserving, this inequality gives that
\begin{align}
   \norm{H^n_s(x)}&\geq \frac{1}{|V_n|}\sum_{v\in V_n}\norm{x+s\psi_n(v)}-\frac{2k}{h}\omega\left(\Lip(\psi_n)\right)
  \label{eq:key}
\end{align}
for all  $n\in\N$, all 
      $x\in B_{X_n}$, and all  $ s\in[0,1]$.
As $ \sum_{v\in V_n}\psi_n(v)=0$ (see \eqref{Eq.prop.psi}), 
\begin{align*}
  \frac{1}{|V_n|}\sum_{v\in V_n}\norm{x+s\psi_n(v)}\geq 
   \left\|x+s\frac{1}{|V_n|}\sum_{v\in V_n}\psi_n(v)\right\|  \geq\left\|x\right\|, 
\end{align*}
and it follows that 
  \begin{equation} \label{Eq.lowerboundH1-eps}          \norm{H^n_s(x)} \geq 1-\frac{2k}{h}\omega\left(\Lip(\psi_n)\right)\ \text{ for all }\ x\in S_{X_n}\ \text{ and }\ s\in [0,1].\end{equation}
      

By our choice of $(\psi_n)_n$,  $\lim_n\Lip(\psi_n)=0$ (see \eqref{Eq.prop.psi}). Hence, since  $\lim_{t\to 0}\omega(t)=0$, we can  pick $n_0\in\N$ such that   \linebreak$4k\omega (\Lip(\psi_n))<h$ for all $n\geq  n_0$. By \eqref{Eq.lowerboundH1-eps},   $\norm{H^n_s(x)} \geq 1/2$ for all $n\geq n_0$, $x\in S_{X_n}$, and $s\in [0,1]$. We can then define,  for any given $n\geq n_0$,
\begin{equation}
  f^n_s(x)=\frac{H^n_s(x)}{\norm{H^n_s(x)}}, \ \text{ for all }\ x\in S_{X_n} \ \text{ and }\ s\in  [0,1].
  \label{eq:homotopy}
\end{equation}
So, $(f^n_s)_{s\in [0,1]}$ is a homotopy starting at $f^n_0=F_n$. This completes the first step in our approach towards applying Proposition \ref{Prop.null.homotopy}. 

We are   left to show that  $f^n_1$ is null-homotopic for $n$ sufficiently large; for this we show that $H^n_1$ is bounded away from zero on $B_{X_n}$. Notice that, using that $\sum_{v\in V_n} \psi_n(v)=0$ once again,
\begin{align*}
  \frac{1}{|V_n|}\sum_{v\in V_n}\norm{\psi_n(v)} & \leq\frac{1}{|V_n|}\sum_{v\in V_n}\norm{x+\psi_n(v)}+\norm{x}\\
  &\leq\frac{1}{|V_n|}\sum_{v\in V_n}\norm{x+\psi_n(v)}+\norm{x+\frac{1}{|V_n|}\sum_{v\in V_n}\psi_n(v)}\\
 &\leq  \frac{2}{|V_n|}\sum_{v\in V_n}\norm{x+\psi_n(v)}.
\end{align*}
Therefore, applying  \eqref{eq:key} with $s=1$, we get that 
\begin{equation}
 \norm{H^n_1(x)}\geq  \frac{1}{2|V_n|}\sum_{v\in V_n}\norm{\psi_n(v)}-\frac{2k}{h}\omega\left(\Lip(\psi_n)\right)\ \text{ for all }\ x\in B_{X_n}.
  \label{eq:ball}
\end{equation}

By our choice of $(\psi_n)_n$ (see \eqref{Eq.prop.psi}), we can pick  $n_1\geq n_0$ such that
\[\frac{1}{2|V_n|}\sum_{v\in V_n}\norm{\psi_n(v)}\geq \frac{1}{4}\ \text{ and }\ \frac{2k}{h}\omega\left(\Lip(\psi_n)\right)\leq\frac{1}{8}\ \text{ for all } \ n\geq n_1.\]
Then  \eqref{eq:ball} gives that 
\begin{equation*}
 \norm{H^n_1(x)}\geq\frac{1}{8}\ \text{ for all }\ x\in B_{X_n}.
\end{equation*}
By Proposition \ref{Prop.null.homotopy}, this shows that $f^n_1$ is null-homotopic for all $n\geq n_1$.\end{proof}


 \begin{AIusage}
Theorem \ref{Thm.main} was obtained by  GPT-6 Astra with little input from the authors.   The organization and exposition of the paper are the authors' responsibility only. Moreover, all work done in collaboration with GPT-6 Astra  was independently checked and reworked by the authors,   who also take full responsibility for the correctness  and integrity of the work. 
\end{AIusage}

 \begin{leanformalization*}
All  mathematical results of this paper were formalized using ChatGPT Astra in Lean 4 (version 4.33.1), using the Mathlib library. While all mathematics developed within  this paper was formalized, we  did not formalize   Osajda's result showing that there are finitely generated groups which contain isometric copies of expander graphs (\cite[Theorem 4]{Osajda2020Acta}). Consequently, Corollary \ref{Cor.no.PropH.cL.universal} is formalized modulo this result. 
 \end{leanformalization*}
 
\begin{disclaimer*}
Other proofs of $c_0$ failing Property (H)
and of Corollary~\ref{Corollary.Johnson} appeared on arXiv on September 15th, 2026 in \cite{ChengChengWang}. We obtained the proofs of the results in this paper by September 13th, 2026, independent of \cite{ChengChengWang}. While results being obtained independently by different groups has always been part of mathematical research, we believe that this may become more common with the advent of new technologies. After some consideration, we decided to still present this paper as our final draft is short and the main results, Theorem~\ref{Thm.main} and Corollary~\ref{Cor.Prop.H}, cover a potentially larger class of Banach spaces than just those of trivial cotype.
\end{disclaimer*}


\begin{acknowledgements*} This paper was written under the auspices of the
American Institute of Mathematics (AIM) SQuaREs program as part of the
``Nonlinear geometry of Banach spaces''  SQuaRE project.
\end{acknowledgements*}


\providecommand{\bysame}{\leavevmode\hbox to3em{\hrulefill}\thinspace}
\providecommand{\MR}{\relax\ifhmode\unskip\space\fi MR }
% \MRhref is called by the amsart/book/proc definition of \MR.
\providecommand{\MRhref}[2]{%
  \href{http://www.ams.org/mathscinet-getitem?mr=#1}{#2}
}
\providecommand{\href}[2]{#2}
\begin{thebibliography}{CCW26}

\bibitem[Aha74]{Aharoni1974Israel}
I.~Aharoni, \emph{Every separable metric space is {L}ipschitz equivalent to a
  subset of {$c\sp{+}\sb{0}$}}, Israel J. Math. \textbf{19} (1974), 284--291.
  \MR{511661}

\bibitem[CCW26]{ChengChengWang}
L.~{Cheng}, Q.~{Cheng}, and Y.~{Wang}, \emph{{On Property (H) of $c_0$ and the
  Sphere Problems of Gromov and Johnson}}, arXiv e-prints (2026),
  arXiv:2609.15817.

\bibitem[CW18]{ChengWang2018JMAA}
Q.~Cheng and Q.~Wang, \emph{On {B}anach spaces with {K}asparov and {Y}u's
  {P}roperty ({H})}, J. Math. Anal. Appl. \textbf{457} (2018), no.~1, 200--213.
  \MR{3702702}

\bibitem[Joh11]{JohnsonOverFlow}
W.~B. Johnson, \emph{Lipschitz homeomorphisms between spheres of n-dimensional
  spaces}, 2011, URL:https://mathoverflow.net/q/60426 (version: 2011-04-03).

\bibitem[KY12]{KasparovYu2012GeoTop}
G.~Kasparov and G.~Yu, \emph{{The Novikov conjecture and geometry of Banach
  spaces}}, Geometry \& Topology \textbf{16} (2012), no.~3, 1859 -- 1880.

\bibitem[Lub10]{LubotzkyBook2010}
A.~Lubotzky, \emph{Discrete groups, expanding graphs and invariant measures},
  Modern Birkh\"{a}user Classics, Birkh\"{a}user Verlag, Basel, 2010, With an
  appendix by Jonathan D. Rogawski, Reprint of the 1994 edition.

\bibitem[OS94]{OdellSchlumprecht1994Acta}
E.~Odell and Th. Schlumprecht, \emph{The distortion problem}, Acta Math.
  \textbf{173} (1994), no.~2, 259--281. \MR{1301394}

\bibitem[Osa20]{Osajda2020Acta}
D.~Osajda, \emph{Small cancellation labellings of some infinite graphs and
  applications}, Acta Math. \textbf{225} (2020), no.~1, 159--191. \MR{4176066}

\bibitem[Ost13]{Ostrovskii2013Book}
M.~Ostrovskii, \emph{Metric embeddings}, De Gruyter Studies in Mathematics,
  vol.~49, De Gruyter, Berlin, 2013, Bilipschitz and coarse embeddings into
  Banach spaces. \MR{3114782}

\end{thebibliography}


\end{document}